\section{反応の速さと平衡を分けて考える}
\label{theory:chem_eq:rates}
同じ反応でも「どちらへどこまで進むか」と「どれほど速く進むか」は別の問いである。前者は反応によるギブズエネルギーの変化、後者は経路上の障壁や反応物の遭遇で決まる。平衡とは分子運動の停止ではなく、正逆の反応の釣合いである。
例えば両方向が一次の \(\mathrm{A\rightleftharpoons B}\) では、
\[
v_f=k_f[A],\qquad v_r=k_r[B],\qquad
v_f=v_r\ \Longrightarrow\ \frac{[B]}{[A]}=\frac{k_f}{k_r}.
\]
速度の比 \(v_r/v_f=1\) と速度定数の比は異なる。一般の複雑な反応では、平衡定数を任意の見かけの速度定数比に置き換えることもできない。
\begin{tcolorbox}[title=原典の不一致を保持する]
令和2年度応用理学13の第3肢は「反応速度」の比を平衡定数とする。原典画像でもこの語を確認した。独立初回は3、公式は2であり、初回を保持している。第2肢の断面積の定義にも確認事項があるため、一意の誤りを説明できた問題として扱わない。「反応速度定数」への無断置換は行わない。
\end{tcolorbox}

\subsection{反応ギブズエネルギーから平衡定数へ}
\label{theory:chem_eq:gibbs}
一定温度・圧力で、反応進行度を \(\xi\)、生成物を正とする化学量論係数を \(\nu_i\) とすると
\[
\Delta_rG=\left(\frac{\partial G}{\partial\xi}\right)_{T,p}
=\sum_i\nu_i\mu_i,\qquad
\mu_i=\mu_i^\circ+RT\ln a_i.
\]
ここで \(a_i\) は標準状態に対する無次元の活量である。これを代入すれば
\[
\Delta_rG=\Delta_rG^\circ+RT\ln Q,\qquad
Q=\prod_i a_i^{\nu_i},\qquad
K=\exp\!\left(-\frac{\Delta_rG^\circ}{RT}\right).
\]
最後の式は平衡の \(\Delta_rG=0\) と \(Q=K\) から得られる。負の \(\Delta_rG\) は正方向への進行が \(G\) を下げることを表す。標準量は指定した標準状態の値であるが、実際の反応量は組成に依存する。標準生成ギブズエネルギーからは、生成物の和と反応物の和の差で \(\Delta_rG^\circ\) を求める。\cite{iupac_greenbook2012}

\subsection{操作直後の反応商を比較する}
\label{theory:chem_eq:quotient}
操作後の \(Q<K\) は生成方向、\(Q>K\) は逆方向への駆動を示す。「圧力が上がれば分子数の少ない側」という略則だけを使わず、何の分圧が変わったかを調べる。
\label{theory:chem_eq:inert}
理想気体では \(p_i=n_iRT/V\) である。定温定積で不活性気体を加えても反応種の分圧は変わらず、全圧だけが上がる。定温で体積を変える操作とは異なる。
\(\mathrm{N_2O_4\rightleftharpoons2NO_2}\) なら、標準圧 \(p^\circ\) を用いた
\[
Q=\frac{(p_{\mathrm{NO_2}}/p^\circ)^2}{p_{\mathrm{N_2O_4}}/p^\circ}
\]
は、組成を固定した瞬間には \(V^{-1}\) に比例する。体積増加は \(Q\) を下げ、解離側へ駆動する。
\label{theory:chem_eq:ammonia}
アンモニア合成 \(\mathrm{N_2+3H_2\rightleftharpoons2NH_3}\) では圧縮が生成側に有利である。発熱反応の温度依存性は
\[
\frac{d\ln K}{dT}=\frac{\Delta_rH^\circ}{RT^2}
\]
で表され、発熱なら昇温で \(K\) は減る。低温は平衡上有利でも速度が遅くなるため、実際の製造では温度・圧力・触媒を組み合わせる。

\subsection{触媒と活性化障壁}
\label{theory:chem_eq:catalyst}
触媒は反応機構に参加し、正味で再生される。「自ら変化しない」は途中で中間体を作らないという意味ではない。別の経路の障壁を下げても、同じ反応物と生成物のエンタルピー差や標準ギブズエネルギー差は変えない。従って反応熱と平衡定数も変えない。選択性は競合する経路の相対速度の差から生じる。
古い問題の「負触媒」は速度を下げる物質を指す呼称であるが、現行IUPACは触媒・触媒作用という語を速度低下に用いず阻害を区別する。用語の時点を分けて読む。\cite{iupac_catalyst}
\label{theory:chem_eq:activation}
アレニウス型の \(k=A\exp(-E_a/RT)\) は温度依存性の模型である。経験的活性化エネルギーの定義は
\[
E_a=RT^2\frac{d\ln k}{dT}.
\]
単純な障壁模型では反応物から遷移状態への差で説明できるが、温度依存の前因子、複数の経路、溶媒などがあると一致しない。電子的な障壁高さと、測定した速度定数の温度勾配から定まる量を区別する。\cite{iupac_greenbook2012,mit_kinetics2013}

\subsection{衝突頻度・断面積・拡散律速}
\label{theory:chem_eq:collision}
分子Aが数密度 \(n_B\) のBに出会う頻度は、単純模型で
\[
z_A(B)=n_B\sigma\overline v_{\mathrm{rel}}\quad[\mathrm{s^{-1}}],\qquad
Z_{AB}=n_Az_A(B)\quad[\mathrm{m^{-3}s^{-1}}].
\]
剛体球なら \(\sigma=\pi(r_A+r_B)^2\) である。温度と組成一定の理想気体で、1個あたりの \(z\) は圧力に比例し、体積あたりの \(Z\) は圧力の2乗に比例する。単位を見ればどちらかを区別できる。
立体因子は単純な衝突模型で反応に適した配向の割合を表す。一般の反応断面積は生成物への流束を入射流束で割った有効面積であり、エネルギーに依存する。幾何学的な剛体球断面積との大小を、比較対象の定義を指定せず一般則にしない。\cite{iupac_greenbook2012}
\label{theory:chem_eq:diffusion}
溶液内で \(\mathrm{A+B\rightleftharpoons AB\to P}\) と遭遇対を経由するとき、ABが分離するより速くPを作る極限では、全体の速さは出会うまでの拡散で制限される。逆に遭遇しても反応しにくければ反応過程が制限する。高速の酸塩基中和を考える際はこの二つを分ける。

\section{酸塩基平衡を組成の問題として解く}
\subsection{酸の定義は受け渡す対象で整理する}
\label{theory:chem_eq:definitions}
アレニウスは水中の水素イオン、ブレンステッドはプロトン、ルイスは電子対に着目する。ブレンステッド酸はプロトン供与体、ルイス酸は電子対受容体である。ルクス・フラッド酸は酸化物イオン受容体、溶媒系の酸は自己解離による陽イオンを増やす物質である。酸の強さと酸化剤としての強さは違う。相対的なプロトン授受では、酢酸も強酸溶媒中で塩基になり得る。
\label{theory:chem_eq:acid_strength}
酸 \(\mathrm{HA}\) の解離は
\[
K_a=\frac{a_{\mathrm{H^+}}a_{\mathrm{A^-}}}{a_{\mathrm{HA}}},
\qquad \mathrm{p}K_a=-\log_{10}K_a.
\]
\(K_a\) が大きいほど、従ってpKaが小さいほど酸は強い。問題の「酸解離定数pKa」は用語として不正確であり、Kaの負対数を指すことを明記する。同じ塩素中心のオキソ酸では、酸素による共役塩基の安定化からHClO2がHClOより強い。\cite{purdue_acid_strength}
HFは通常の水溶液で弱酸に分類されるが、それは有害性が小さいという意味ではない。ガラスとの反応や高濃度溶液の発煙は別の性質である。\cite{niosh_hf}

\subsection{中和を先に、加水分解を後に}
\label{theory:chem_eq:hydrolysis}
強酸・強塩基による量的な中和を物質量で処理し、その後に残った化学種の平衡を解く。弱酸と強塩基の等モル中和では、過剰水酸化物イオンがなくても、共役塩基の加水分解で塩基性になる。
\[
\mathrm{A^-+H_2O\rightleftharpoons HA+OH^-},\qquad K_b=K_w/K_a.
\]
濃度近似で塩の形式濃度を \(C\)、加水分解量を \(x\) とすれば
\[
K_b=\frac{x^2}{C-x},\qquad x\simeq\sqrt{K_bC}\quad(x\ll C).
\]
混合体積を分母にして \(C\) を求める。平成30年度の問題では両液の各0.020 molを中和し、全体積0.020 dm3から \(C=1\,\mathrm{mol\,dm^{-3}}\) となる。与えられた定数でpHは9となる。この計算は活量係数を1とする近似であり、濃厚溶液の実測値を保証する式ではない。水の定数の温度依存性もあるので、pHとpOHの和を常に14と覚えず与えた \(K_w\) を使う。

\subsection{共役対と緩衝容量}
\label{theory:chem_eq:buffer}
弱酸と共役塩基を共存させると、加えた塩基を酸が、加えた酸を共役塩基が消費する。酸解離式を対数化すれば
\[
\mathrm{pH=p}K_a+\log_{10}\frac{a_{\mathrm{A^-}}}{a_{\mathrm{HA}}}
\]
を得る。濃度比への置換や調製時の物質量比への置換には、活量係数や解離による濃度変化の近似が必要である。
比がpHを決め、全量がどれだけの添加に耐えるかを左右する。片方を使い切れば緩衝は失われる。弱塩基と強酸を混ぜる場合も、部分中和して弱塩基と共役酸を共存させる必要がある。血液のpH維持にも複数の酸塩基対による緩衝が関与する。\cite{mit_buffer2008}

\subsection{溶解度積とpH}
\label{theory:chem_eq:solubility}
\(\mathrm{Mg(OH)_2(s)\rightleftharpoons Mg^{2+}+2OH^-}\) では、固体の活量を1として
\[
K_{\mathrm{sp}}=a_{\mathrm{Mg^{2+}}}a_{\mathrm{OH^-}}^2
\]
となる。一定温度で定数が一定でも、pHによって水酸化物イオンの活量が変われば溶解量が変わる。純水での溶解だけから得た \([\mathrm{OH^-}]=2s\) を、外部から酸や塩基を加えた系へそのまま適用しない。

\subsection{出題の比較方法}
アンモニア合成の平成28・30・令和4年度は触媒の肢番号が入れ替わる。令和3・5年度のギブズエネルギーも同じ誤りが別番号にある。番号を覚えるより、反応商・標準量・速度のどれを問うかを先に見分ける。類似問題の関係と概念頻度は、全問の校正が済むまで暫定範囲の集計として読む。
